% ! TeX root = ../main.tex

\section{Introduction}


The \emph{bulk-synchronous parallel} (BSP)~\cite{Valia1990:BSP} paradigm powers critical
computational workloads, from scientific simulations to large-scale machine learning training.
These systems are uniquely fragile: at the scale of tens of thousands of processors, lockstep
execution means a single straggler can block cluster-wide progress at every synchronization point.
As ML training pushes BSP toward 100K GPUs~\cite{Salpe2026:100K} and scientific workloads grow in
tandem~\cite{Das2023:LargeScale,Stock2024:LargeScale}, the cost of undiagnosed performance
pathologies has become economically untenable.
Yet the causes are diverse and highly
context-dependent~\cite{\citeGroupVarHpc,\citeGroupVarAmr,\citeGroupVarRdma,\citeGroupVarMisc,\citeGroupML}
---load imbalance, poor task overlap, hardware variability---varying across clusters, workloads, and
scales.
Metrics~\cite{:Prometheus,Schwa2021:LDMS} provide baseline visibility but rarely enough workload
structure to attribute causes.
Tracing~\cite{Shend2006:TAU,Knupf2012:Score-P,2025:PytorchKineto,Devar2024:DFTracer,Boehm2016:Caliper}
can capture the fine-grained patterns needed for diagnosis, but prevailing approaches---exhaustive
collection to disk followed by expensive \emph{extract-transform-load}
(ETL)~\cite{Yildi2023:IOMax,\citeJainAmr} ---cannot provide feedback at the speed problems demand.
In a prior study of placement in AMR codes~\cite{\citeJainAmr}, isolating the root causes of
several performance artifacts took months, while fixes typically took days.
Rapid visibility into production pathologies remains one of the highest-yield paths to efficiency.

We argue for always-on, steerable observability as a production capability for BSP workloads, not
just faster post-hoc analysis, but continuous visibility with the ability to escalate on demand.
The evolution of distributed tracing~\cite{Barha2004:Magpie,Fonse2007:X-Trace,Sigel2010:Dapper}
offers an instructive parallel.
In microservices, tracing evolved to preserve causal relationships between requests (the unit of
work in asynchronous systems) so tail latency is continuously monitored and attributable to
specific services.
In BSP, causality is induced by synchronization: performance anomalies manifest as stragglers that
stall collectives.
The same synchronization creates natural windows for both observation and intervention, enabling
operators to steer collection on demand, test hypotheses live, and act on feedback while the
job runs.
Real-time diagnostics are already proving valuable in ML training~\cite{\citeGroupML} for problems
ranging from straggler localization to training correctness, but each system implements a bespoke
collection and analytics pipeline.
% And because BSP is a single coordinated workload, the system can unify observation and% intervention---enabling operators to steer collection on demand, test hypotheses live, and act on% feedback, all while the job runs.
%A common observability substrate is needed for BSP workloads, 
BSP workloads need a common observability substrate, providing low-overhead on-fabric transport,
BSP-aware coordination semantics for telemetry aggregation and control, and programmable analytics
interfaces (\S\ref{sec:orca:req}).

We present a tree-based overlay~\cite{Arnol2006:TBONs} for rapid feedback loops over large-scale BSP
applications called 
ORCA~(\emph{Observability~with~Real-time~Control~and~Aggregation}).
ORCA treats telemetry as columnar streams and provides in-situ SQL-based
analytics, faster offline analytics, and a control plane for
online steerability.
Central to the design is treating synchronization as the fundamental consistency boundary.
On the data path, this yields \emph{timestep dataframes}~(\S\ref{sec:orca:des_model}): a
per-timestep grouping of all ranks' telemetry into a single logical analysis unit.
Because synchronization partitions BSP execution into causally independent windows, timestep
dataframes support both in-situ analysis during execution and timestep-partitioned persistence
for faster offline queries.
Combined with columnar formats (Arrow~\cite{:ApacheArrow} in transit, Parquet~\cite{:ApacheParquet} at rest), this eliminates the ETL
that dominates conventional tracing workflows.
An RDMA-based transport~(\S\ref{sec:orca:des_xport}) minimizes collection overhead,
while distributed query planning with operator pushdown minimizes data movement~(\S\ref{sec:orca:des_insitu}),
On the control path, a timestep-consistent commit protocol (\S\ref{sec:orca:des_ctl}) enables atomic
interventions---new analytics, probes, or parameter updates---applied asynchronously at the same timestep
across all ranks.
% Together, these enable workflows ranging from faster offline analytics to low-overhead, real-time, closed-loop feedback.
We evaluate ORCA on a real AMR code at 512--4096 CPU ranks:

% control paths to realize low-overhead, real-time, closed-loop feedback.
% We evaluate ORCA on a real AMR code at 512--4096 CPU ranks:

% We present a tree-based overlay for real-time observability and control of large-scale BSP applications called ORCA.
% ORCA treats telemetry as columnar streams and provides capabilities ranging from in-situ SQL-based analytics,
% faster dataframe-style offline analytics, and a control plane for online steerability.
% % ORCA provides timestep-consistent control and data
% ORCA constructs a tree-based overlay providing data and control paths
% both treating synchronization as the fundamental consistency boundary in BSP execution.
% On the data path, this yields an abstraction called timestep dataframe, consisting of all ranks' telemetry from a single timestep.
% As timesteps are causally independent, timestep dataframes enable in-situ analyses, and timestep-partitioned 
% layouts for faster offline analyses. Combined with columnar formats like Arrow and Parquet, this enables

% ORCA treats the synchronization window as the fundamental consistency boundary in BSP execution

% We present a \emph{tree-based overlay network} (TBON)~\cite{Arnol2006:TBONs} for real-time
% observability and control of large-scale BSP applications, called
% ORCA~(\emph{Observability~with~Real-time~Control~and~Aggregation}).
% ORCA treats telemetry as columnar streams, transformed in-situ via SQL-based analytics 
% %interface called \emph{OrcaFlow} 
% en route to \emph{sinks} such as dashboards or parallel filesystems.



% A key contribution is the timestep dataframe abstraction: recognizing that global synchronization
% divides BSP execution into causally independent windows, ORCA makes timestep the key thing
% for in-situ and offline  something.
% Timestep dataframes, consisting of all telemetry from a single timestep, enable both faster offline
% analysis by partitioning, and in-situ analysis too.
% Distributed query planning with operator pushdown minimizes data movement for in-situ analysis,
% (\S\ref{sec:orca:des_insitu}), while an end-to-end columnar pipeline enables zero-ETL post-mortems.

% A key contribution is the \emph{timestep dataframe} abstraction: building on the insight that
% global synchronization divides BSP telemetry into causally independent windows, ORCA transforms
% telemetry from a rank-major to a timestep-major representation (\S\ref{sec:orca:des_model}).
% An RDMA-based TBON-optimized transport (\S\ref{sec:orca:des_xport}) enables low-overhead collection of
% these dataframes.
% They can then be transformed via SQL-based analytics, routed to a dashboard, or persisted as
% partitioned Parquet for zero-ETL post-mortems.
% Distributed query planning with operator pushdown minimizes data movement across the overlay
% (\S\ref{sec:orca:des_insitu}).

% A control plane (\S\ref{sec:orca:des_ctl}) provides atomic and timestep-consistent control semantics,
% guaranteeing correctness for interventions such as new analytics, probes, parameter updates.
% ORCA combines a modern columnar analytics
% stack~\cite{:ApacheArrow,:ApacheParquet,Lamb2024:DataFusion} with RDMA and BSP-aligned data and
% control paths to realize low-overhead, real-time, closed-loop feedback.
% We evaluate ORCA on a real AMR code at 512--4096 CPU ranks:

\begin{itemize}[leftmargin=*]
  \item In conventional tracing workflows, ORCA captures equivalent telemetry at \mytilde{2}\% runtime overhead
        (vs 56--81\% for state-of-the-art tracers), produces up to 42$\times$ smaller traces, and
        enables up to 7275$\times$ faster offline analyses (\S\ref{sec:orca:eval-trace}).
  \item In-situ analytics enable low-overhead needle-in-the-haystack
        workflows---outlier \mpiw{} calls are isolated in real-time at sub-2\% runtime
        overhead, while discarding up to 99.999\% of data at MPI ranks (\S\ref{sec:orca:eval-insitu}).
  \item In an agentic evaluation of its steerable capabilities using an LLM
        agent,
        ORCA enables localizing a performance anomaly---the same bug from \cite{\citeJainAmr} that took
        months of manual investigation---in 27 minutes, while reducing data volume by up to 144$\times$
        (\S\ref{sec:orca:eval-agent}).
\end{itemize}

More broadly, ORCA demonstrates that declarative SQL analytics can be embedded on the fabric
alongside a live MPI application, enabling streaming analytics workflows beyond observability~(\S\ref{sec:orca:related}).
ORCA is open source~\cite{Autho2026:orca-code}.
